\chapter{模板使用指南}

\section{修改个人信息}
你需要打开主文件 \texttt{main.tex}，在导言区（\texttt{\textbackslash begin\{document\}} 之前）修改以下命令：
\begin{itemize}
    \item \texttt{\textbackslash title\{...\}}：论文中文题目。
    \item \texttt{\textbackslash author\{...\}}：你的姓名。
    \item \texttt{\textbackslash studentid\{...\}}：学号。
    \item \texttt{\textbackslash college\{...\}}：学院（如理学院）。
    \item \texttt{\textbackslash major\{...\}}：专业。
    \item \texttt{\textbackslash supervisor\{...\}}：指导教师姓名。
\end{itemize}

\section{章节与段落写作建议}

\subsection{如何增加新章节}
建议在 \texttt{contents/} 目录下新建 \texttt{.tex} 文件，例如 \texttt{05-new-chapter.tex}。然后在 \texttt{main.tex} 中使用 \texttt{\textbackslash include\{contents/05-new-chapter\}} 引入。

\subsection{段落写作规范}
\begin{enumerate}
    \item \textbf{首行缩进}：中文会首行自动缩进。你只需要正常撰写，换行时空一行即可开始新段落。
    \item \textbf{强调文本}：使用 \texttt{\textbackslash textbf\{粗体\}} 或 \texttt{\textbackslash textit\{斜体\}}。
    \item \textbf{列表环境}：
        \begin{itemize}
            \item \texttt{itemize} 用于无序列表（如本列表）。
            \item \texttt{enumerate} 用于有序列表（自动生成 1, 2, 3）。
        \end{itemize}
\end{enumerate}

\section{Overleaf 编译超时解决方法}
若在 Overleaf 中遇到编译超时（Timeout），可尝试以下方案：
\begin{enumerate}
    \item \textbf{局部编译}：在 \texttt{main.tex} 导言区使用 \texttt{\textbackslash includeonly\{contents/正在写的章节\}}。
    \item \textbf{草稿模式}：将 \texttt{\textbackslash documentclass} 改为 \texttt{\textbackslash documentclass[draft]\{muc-thesis\}}，可跳过图片渲染。
    \item \textbf{优化图片}：确保 \texttt{figures/} 下的图片单张大小不超过 1MB，推荐使用 PDF 或压缩后的 PNG。
\end{enumerate}

\remark{
\noindent 关于其他具体内容相关的问题，请查看word模板中的相关说明，或者直接联系指导教师进行咨询。

\noindent 注意，本模板最初学生自主开发，虽然目前我们已经有多位老师和同学协力测试和优化，但仍可能存在一些细节问题或不适用的情况。因此我们强烈建议，即使使用了本模板，也要仔细阅读学校的相关规范文件，\textbf{同时也要阅读word模板中的相关说明}，确保最终提交的论文符合学校的要求。

\noindent 如果你在使用过程中遇到任何问题，欢迎随时反馈。请认准项目官方Github链接：\url{https://github.com/KardeniaPoyu/muc-thesis}，在Issues中提交你的问题或建议，我们会尽快回复并持续优化模板。
}